\documentclass{article}
\usepackage{amsmath,amssymb}
\usepackage{xurl}
\usepackage[colorlinks=true,linkcolor=blue,citecolor=blue,urlcolor=blue]{hyperref}
% Shared commands for notes in docs/paper-gaps/.

\DeclareUnicodeCharacter{2080}{\ifmmode{}_{0}\else\textsubscript{0}\fi}
\DeclareUnicodeCharacter{2081}{\ifmmode{}_{1}\else\textsubscript{1}\fi}
\DeclareUnicodeCharacter{2082}{\ifmmode{}_{2}\else\textsubscript{2}\fi}
\DeclareUnicodeCharacter{2099}{\ifmmode{}_{n}\else\textsubscript{n}\fi}

\newcommand{\Lean}{\textsc{Lean}}
\ExplSyntaxOn
\NewDocumentCommand{\leanid}{m}{
  \ifmmode
    \textsf{\detokenize{#1}}
  \else
    \group_begin:
    \urlstyle{sf}
    \tl_set:Nn \l_tmpa_tl {#1}
    \tl_replace_all:Nnn \l_tmpa_tl {\_} {_}
    \exp_args:NV \path \l_tmpa_tl
    \group_end:
  \fi
}
\ExplSyntaxOff
\providecommand{\tr}{\operatorname{tr}}
\newcommand{\tnleanIssueBase}{https://github.com/LionSR/TNLean/issues/}
\newcommand{\tnleanPrBase}{https://github.com/LionSR/TNLean/pull/}
\newcommand{\tnleanDocBase}{https://sirui-lu.com/TNLean/docs/}
\newcommand{\ghissue}[1]{\href{\tnleanIssueBase#1}{GitHub issue \##1}}
\newcommand{\ghpr}[1]{\href{\tnleanPrBase#1}{PR \##1}}
\newcommand{\doclink}[1]{\href{\tnleanDocBase#1.html}{\path{#1}}}

% Dirac notation and the identity operator, as in blueprint/src/macros/common.tex.
% \providecommand keeps note-local definitions authoritative.
\usepackage{dsfont}
\providecommand{\ket}[1]{|#1\rangle}
\providecommand{\bra}[1]{\langle #1|}
\providecommand{\braket}[2]{\langle #1 | #2 \rangle}
\providecommand{\ketbra}[2]{|#1\rangle\!\langle #2|}
\providecommand{\Id}{\mathds{1}}
\providecommand{\one}{\mathds{1}}
\providecommand{\C}{\mathbb{C}}
\providecommand{\MN}[1]{M_{#1}(\C)}
\providecommand{\MD}{M_D(\C)}

% Verdict marker: \gapnote{<kind>}{<status>}. Kind is the policy
% classification (clarification, local-correction, scope-restriction,
% unfaithful, false-source, open-gap); status is open, wip (elimination
% underway), resolved, or historical. Machine-read by the site index;
% prints a verdict line.
\newcommand{\gapnote}[2]{\par\noindent\textbf{Verdict:} #1 (#2).\par}

\title{Connected Domains in the One-Dimensional Area Law}
\author{}
\date{October 16, 2026}
\begin{document}
\maketitle
\gapnote{scope-restriction}{open}

\section{The source statement}

Theorem~A of the one-dimensional note~\cite{TNLean2026OneDimensional} is
stated for every finite set $\Lambda\subset\mathbb Z$ with its induced path
graph: for an admissible Hamiltonian $H=\sum_X h_X$ on $\Lambda$, one Hermitian
term of norm at most $J$ per nonempty support $X\subseteq\Lambda$ of graph
diameter at most $R$, and a unit ground vector $\Omega$ with
$H-E_0\ge\Delta(1-|\Omega\rangle\langle\Omega|)$, every $A\subseteq\Lambda$
satisfies $S_\Omega(A)\le C_{\mathrm{AL}}\,|\partial_\Lambda A|$ with
$C_{\mathrm{AL}}\le C_0(q,R)(1+J/\Delta)$. The edge boundary
$\partial_\Lambda A$ counts the edges $\{x,x+1\}\subseteq\Lambda$ with exactly
one endpoint in $A$.

The note's first reduction (Section~2) disposes of disconnected domains in a
sentence: the Hamiltonian is a sum over the components of $\Lambda$, each
component is an interval of consecutive integers, the unique gapped ground
vector is a product of the component ground vectors, and both the entropy and
the edge boundary add over components. The theorem is then proved for
connected $\Lambda$, that is, for a chain of $L$ consecutive sites.

\section{What is proved and the restriction}

The formalization states the model on the chain of $L$ sites, with sites
$\{0,\ldots,L-1\}$, supports of range $R$ the nonempty sets of sites of
diameter at most $R$, and the Hamiltonian, gapped ground vector, regional
entropies and boundary edges as in the note.\footnote{Declarations:
\leanid{TNLean.MPS.AreaLaw.IsChainSupport},
\leanid{TNLean.MPS.AreaLaw.ChainHamiltonian},
\leanid{TNLean.MPS.AreaLaw.boundaryEdges}.} Theorem~A is proved on this
chain for every subset $A$, with the explicit constant
$C_0(q,R)(1+J/\Delta)$.\footnote{Declarations:
\leanid{TNLean.MPS.AreaLaw.regionEntropy_interval_le},
\leanid{TNLean.MPS.AreaLaw.regionEntropy_le_mul_card_boundaryEdges},
\leanid{TNLean.MPS.AreaLaw.chainAreaLaw}.} Theorem~B of the note is stated on
the chain of $L$ sites in the note itself, so it carries no restriction.

The restriction is to connected domains. For a disconnected
$\Lambda\subset\mathbb Z$ the interval bound still holds verbatim on every
interval of positions, since the proof uses only the crossing count and the
buffered discount; but the subset bound $S_\Omega(A)\le C|\partial_\Lambda A|$
needs the product structure of the ground vector over components (a region
inside one component contributes no boundary edge at a gap of $\Lambda$, so
its entropy must vanish for the bound to hold), and that product structure is
not formalized.

\section{Elimination plan}

Formalize the component reduction: for a Hamiltonian that is a sum of terms
supported in disjoint blocks of sites, a unit gapped ground vector of the sum
is a tensor product of gapped ground vectors of the blocks (the ground energy
adds, and uniqueness of the global ground vector forces the product form), and
the regional entropy of a product vector adds over the blocks. Theorem~A for an
arbitrary finite $\Lambda\subset\mathbb Z$ then follows from the chain version
applied to each component. Until then the restriction is recorded on the
module docstring of \texttt{TNLean/MPS/AreaLaw/AreaLaw.lean} and in the blueprint
statement of Theorem~A.

\bibliographystyle{alpha}
\bibliography{references}
\end{document}
